\section{Particle Filters as Approximate Inference}\label{sec:particle-filters}

\subsection{Sequential importance sampling}

Importance sampling estimates expectations under a target density $p$ from samples of a proposal density $q$ that
is easier to sample: with $\x^{(i)} \sim q$ and weights $\w^{(i)} \propto p(\x^{(i)}) / q(\x^{(i)})$ normalized to sum
to one, $\mathbb{E}_p[f] \approx \sum_i \w^{(i)} f(\x^{(i)})$. The particle filter applies this to the posterior over
trajectories, $p(\x_{0:t} \mid \z_{1:t}, \uu_{1:t})$, and builds the samples one time step at a time
\cite{arulampalam2002tutorial,doucet2011tutorial}. Each of the $N$ particles extends its trajectory with a new pose
drawn from a proposal $q(\x_t \mid \x_{t-1}, \z_t, \uu_t)$, and by the factorization \eqref{eq:joint} its weight
is updated recursively:
\begin{equation}
  \w_t^{(i)} \;\propto\; \w_{t-1}^{(i)}\,
  \frac{p(\z_t \mid \x_t^{(i)})\; p(\x_t^{(i)} \mid \x_{t-1}^{(i)}, \uu_t)}
       {q(\x_t^{(i)} \mid \x_{t-1}^{(i)}, \z_t, \uu_t)}.
  \label{eq:weight}
\end{equation}
The weighted particles approximate the filtering belief \eqref{eq:bayes-filter}, and the pose estimate is their
weighted mean.

The variance of the weights can only grow over time \cite{doucet2011tutorial}. After a few steps almost all the
weight sits on a few particles and the rest contribute nothing, which is called weight degeneracy. A common measure
is the effective sample size
\begin{equation}
  \mathrm{ESS} = \frac{1}{\sum_{i=1}^{N} \bigl(\w_t^{(i)}\bigr)^2},
  \label{eq:ess}
\end{equation}
which is $N$ for equal weights and 1 when one particle holds all the weight. Two things work against degeneracy: a
proposal that places particles where the posterior is high, and resampling.

\subsection{Proposal distributions}

The \emph{motion-model} (bootstrap) proposal $q = p(\x_t \mid \x_{t-1}, \uu_t)$ \cite{gordon1993novel} cancels the
motion term in \eqref{eq:weight}, so the weight is just the likelihood of the observation. It is simple, but it
ignores $\z_t$ when sampling: if the sensor is more precise than the motion, most particles land where the likelihood
is small.

Given $\x_{t-1}$, the proposal that minimizes the variance of the weight is $p(\x_t \mid \x_{t-1}, \z_t, \uu_t)$,
the pose distribution after both the motion and the new observation \cite{doucet2011tutorial}. It rarely has a closed
form, and the measurement-informed proposals approximate it:
\begin{itemize}
  \item The \emph{extended Kalman particle filter} gives each particle an extended Kalman filter. The filter predicts
        the pose with the motion model, updates it with the linearized measurement model, and the new pose is drawn
        from the resulting Gaussian; the weight then follows \eqref{eq:weight} with that Gaussian as $q$
        \cite{elfring2021particle}.
  \item The \emph{auxiliary particle filter} \cite{pitt1999filtering} first decides which particles to continue.
        It propagates each particle once, weights it by how well the propagated pose explains $\z_t$, and resamples
        by these first-stage weights. The selected particles are then propagated again with the motion model, and
        their weights are divided by the first-stage likelihood.
  \item \emph{FastSLAM 2.0} \cite{montemerlo2003fastslam2} samples the pose from a Gaussian that combines the
        predicted pose with the observations of already mapped landmarks, including the uncertainty of those
        landmarks. The weight is the predictive likelihood of the observations. FastSLAM 1.0
        \cite{montemerlo2002fastslam} uses the motion-model proposal.
\end{itemize}
Measurement-informed proposals cost more per particle. They should therefore reach a given accuracy with fewer
particles; whether they pay off at a given runtime is an empirical question.

\subsection{Resampling}

Resampling draws $N$ new particles from the current ones, with probability proportional to their weights, and resets
all weights to $1/N$. Particles with negligible weight disappear and heavy particles are copied. The schemes differ
in how much randomness they add to the number of copies $N_i$ of particle $i$, whose expectation is $N \w^{(i)}$ in
all four \cite{douc2005comparison,hol2006resampling}:
\begin{itemize}
  \item \emph{Multinomial}: $N$ independent draws from the weights, so $N_i$ is binomial.
  \item \emph{Residual}: $\lfloor N \w^{(i)} \rfloor$ copies of each particle deterministically, and the remaining
        particles by multinomial sampling from the residual weights.
  \item \emph{Stratified}: the interval $[0, 1)$ is split into $N$ equal strata, one uniform number is drawn in each,
        and each number selects the particle whose cumulative weight covers it.
  \item \emph{Systematic}: a single uniform number $u \in [0, 1/N)$ gives the $N$ positions $u + (j-1)/N$, which
        select particles in the same way.
\end{itemize}
Residual, stratified and systematic resampling have lower variance than multinomial resampling
\cite{douc2005comparison}, and all four can be implemented in $O(N)$ time.

Resampling removes weight degeneracy but causes sample impoverishment: after resampling many particles are copies of
the same pose, and the diversity lost cannot be recovered. Resampling at every step, as the original bootstrap filter
does, therefore throws away information when the weights are still balanced. Adaptive rules resample only when the
weights have degenerated. The rules compared here resample when $\mathrm{ESS} < \tau N$, or when the largest weight is
too large, $1 / \max_i \w^{(i)} < \tau N$, for a threshold $\tau \in (0, 1)$ \cite{elfring2021particle}.

\subsection{MCMC moves}

Resample-move algorithms \cite{gilks2001following} restore diversity after resampling by applying a Markov chain Monte
Carlo kernel that leaves the posterior invariant, such as a few Metropolis-Hastings steps, to each particle. Copies
of the same particle then move apart without changing the distribution they represent. The experiments in this
report do not use moves; they are a natural next step (\cref{sec:discussion}).
